\documentclass[11pt]{amsart}
\usepackage[margin=1.15in]{geometry}
\usepackage{amsmath,amssymb,amsthm,mathtools,booktabs}
\usepackage{xcolor}
\usepackage[colorlinks=true,linkcolor=blue!50!black,citecolor=blue!50!black,urlcolor=blue!50!black]{hyperref}
\usepackage{enumitem}

\newtheorem{theorem}{Theorem}[section]
\newtheorem{proposition}[theorem]{Proposition}
\newtheorem{lemma}[theorem]{Lemma}
\newtheorem{corollary}[theorem]{Corollary}
\theoremstyle{remark}
\newtheorem{remark}[theorem]{Remark}
\newtheorem{example}[theorem]{Example}

\newcommand{\E}{\mathbb{E}}
\renewcommand{\Pr}{\mathbb{P}}
\newcommand{\N}{\mathbb{N}}
\newcommand{\SN}{S_N}
\newcommand{\cont}{\operatorname{cont}}
\DeclareMathOperator{\cyc}{cyc}

\title[Same-cycle indicator under random transpositions]{The spectral measure of the same-cycle indicator\\ under random transpositions}
\author{Michael K.\ Kinyon}
\address{Department of Mathematics, University of Denver}
\email{mkinyon@du.edu}
\date{\today}

\begin{document}

\begin{abstract}
Let $u\neq v$ be two fixed points of $[N]$ and let $\varepsilon(\sigma)=2\cdot\mathbf 1\{u\text{ and }v\text{ lie in a common cycle of }\sigma\}-1$ on the symmetric group $S_N$.  We show that the spectral measure of $\varepsilon$ with respect to the random-transposition walk has an explicit closed form: it is supported on the partitions $\lambda=(a,b,1^j)$ with $a\ge b\ge1$, and there
\[
\|\Pi_\lambda\varepsilon\|^2=\frac{2(a-b+1)^2}{(a+j+1)^2(b+j)^2}\,\bigl(1-r_\lambda\bigr),
\]
where $r_\lambda$ is the normalised character of a transposition.  The proof is a few lines: $\mathbf 1\{u\sim_\sigma v\}$ is half the increment of the number of cycles under right multiplication by $(u\,v)$, and the Fourier coefficients of the cycle-count function are read off the content polynomial.  Consequences include the exact two-point correlation $\langle\varepsilon,P^m\varepsilon\rangle$ of the same-cycle status along the walk, the bound $|\langle\varepsilon,P^m\varepsilon\rangle|\le 12/m+o(m^{-3})$ uniformly in $N$, the limiting spectral density $4(\alpha-\beta)^2/((1-\alpha)(1-\beta))\,d\alpha\,d\beta$, and the probability $(N+1)(N+2)/(6N(N-1))$ that one random transposition changes the same-cycle status of two points of a uniformly random permutation.  The same spectral measure governs every conjugation-invariant random walk.
\end{abstract}

\maketitle

\section{Introduction and main result}

Let $N\ge2$ and let $\sigma_0,\sigma_1,\sigma_2,\dots$ be the random-transposition walk on $S_N$: $\sigma_{t+1}=\sigma_t\tau_{t+1}$ with $\tau_1,\tau_2,\dots$ independent uniform transpositions.  Fix two points $u\ne v$ of $[N]$ and write $u\sim_\sigma v$ if $u$ and $v$ lie in a common cycle of $\sigma$.  The evolution of the pair $(u,v)$ under the walk is the simplest instance of the coagulation--fragmentation dynamics of the cycle structure studied by Schramm \cite{Sch05}, Berestycki--Durrett \cite{BD06}, Berestycki--Schramm--Zeitouni \cite{BSZ11} and many others: a transposition whose two points lie in one cycle splits it, and merges two cycles otherwise, and the status of $(u,v)$ changes exactly when the transposition separates them or joins their cycles.  Under the uniform measure $\Pr[u\sim v]=\tfrac12$, so the natural two-point quantity is the autocorrelation of the sign
\[
\varepsilon(\sigma)=2\cdot\mathbf 1\{u\sim_\sigma v\}-1\in\{\pm1\}
\]
along the stationary walk, $\E[\varepsilon(\sigma_0)\varepsilon(\sigma_m)]$ with $\sigma_0$ uniform.

The random-transposition walk is driven by a central element of the group algebra, so it is diagonalised by the irreducible characters of $\SN$, with eigenvalue on the $\lambda$-isotypic component
\begin{equation}\label{eq:r}
r_\lambda=\frac{\chi_\lambda(\tau)}{d_\lambda}=\frac{2}{N(N-1)}\sum_{\square\in\lambda}\cont(\square),
\end{equation}
where $\cont(i,j)=j-i$ is the content of the box in row $i$ and column $j$ \cite{DS81}.  Consequently every autocorrelation of a function $f$ on $\SN$ is a moment of its spectral measure $\nu_f(\lambda)=\|\Pi_\lambda f\|^2$:
\[
\langle f,P^mf\rangle=\sum_{\lambda\vdash N}r_\lambda^{\,m}\,\nu_f(\lambda),
\qquad
\|P^df\|^2=\sum_{\lambda\vdash N}r_\lambda^{\,2d}\,\nu_f(\lambda).
\]
Here $L^2(\SN)$ carries the uniform measure, $(Pf)(\sigma)=\binom N2^{-1}\sum_\tau f(\sigma\tau)$, and $\Pi_\lambda=\frac{d_\lambda}{N!}\sum_\rho\chi_\lambda(\rho)R(\rho)$ with $(R(\rho)f)(\sigma)=f(\sigma\rho)$ is the projection onto the $\lambda$-isotypic component of the right-regular representation.  For class functions the spectral measure is simply the square of the Fourier coefficient; the function $\varepsilon$ is not a class function, and one does not expect its spectral measure to be computable.  It is.

\begin{theorem}\label{thm:main}
$\nu_\varepsilon(\lambda)\ne0$ only for $\lambda=(a,b,1^j)$ with $a\ge b\ge1$ and $a+b+j=N$, and for these shapes
\begin{equation}\label{eq:main}
\nu_\varepsilon(\lambda)=\|\Pi_\lambda\varepsilon\|^2=2\,c_\lambda^{\,2}\,(1-r_\lambda),
\qquad
c_\lambda=\frac{a-b+1}{(a+j+1)(b+j)},
\qquad
r_\lambda=\frac{a(a-1)+b(b-3)-j(j+3)}{N(N-1)} .
\end{equation}
Here $c_\lambda=|\langle c,\chi_\lambda\rangle|$ is the modulus of the Fourier coefficient of the number-of-cycles function $c(\sigma)$.
\end{theorem}

Thus, for $\sigma_0$ uniform,
\begin{equation}\label{eq:corr}
\E\bigl[\varepsilon(\sigma_0)\varepsilon(\sigma_m)\bigr]
=\sum_{\substack{a\ge b\ge1\\ a+b+j=N}}\frac{2(a-b+1)^2}{(a+j+1)^2(b+j)^2}\,(1-r_\lambda)\,r_\lambda^{\,m},
\end{equation}
and $\sum_\lambda\nu_\varepsilon(\lambda)=\|\varepsilon\|^2=1$.  Table~\ref{tab:values} gives some values.  Section~\ref{sec:proof} proves the theorem, Section~\ref{sec:cons} draws the consequences announced in the abstract, and Section~\ref{sec:gen} records the generalisations (other conjugation-invariant walks, non-stationary starts) that the proof gives for free.  The computation was motivated by a problem on random Latin squares, where the status of two marked rows in the cycle structure of a column-pair permutation is randomised by switching a pool of intercalates \cite{KPS25}; we say a word about this in Remark~\ref{rem:latin}.

\begin{table}[h]
\centering
\begin{tabular}{r@{\quad}cccccccc}
\toprule
$m$ & 1 & 2 & 4 & 8 & 16 & 64 & 256 & 1024\\
\midrule
$N=4000$ & $0.6663$ & $0.5107$ & $0.3418$ & $0.2037$ & $0.1121$ & $0.02992$ & $0.00727$ & $0.00149$\\
$N=\infty$ & $2/3$ & $23/45$ & $0.3422$ & $0.2042$ & $0.1125$ & $0.03042$ & $0.00776$ & $0.00195$\\
$m\cdot(N=\infty)$ & $0.667$ & $1.022$ & $1.369$ & $1.634$ & $1.801$ & $1.947$ & $1.986$ & $1.997$\\
\bottomrule
\end{tabular}
\medskip
\caption{The autocorrelation $\E[\varepsilon(\sigma_0)\varepsilon(\sigma_m)]$ of the same-cycle sign along the stationary random-transposition walk, from \eqref{eq:corr} and from the limit \eqref{eq:cont}.}
\label{tab:values}
\end{table}

\section{Proof of Theorem~\ref{thm:main}}\label{sec:proof}

\subsection*{Step 1: the indicator is a difference of the cycle count}
Let $c(\sigma)$ be the number of cycles of $\sigma$.  Right multiplication by the transposition $(u\,v)$ sends $u\mapsto\sigma(v)$ and $v\mapsto\sigma(u)$ and fixes every other image.  If $u$ and $v$ lie on one cycle $u\to\sigma(u)\to\dots\to v\to\sigma(v)\to\dots\to u$, the new permutation closes the two cycles $u\to\sigma(v)\to\dots\to u$ and $v\to\sigma(u)\to\dots\to v$; if they lie on different cycles, these are spliced into one.  Hence $c(\sigma(u\,v))-c(\sigma)=+1$ if $u\sim_\sigma v$ and $-1$ otherwise, i.e.
\begin{equation}\label{eq:step1}
\mathbf 1\{u\sim_\sigma v\}=\tfrac12\bigl(1+c(\sigma(u\,v))-c(\sigma)\bigr),
\qquad
\varepsilon=\bigl(R((u\,v))-I\bigr)c .
\end{equation}

\subsection*{Step 2: the Fourier coefficients of the cycle count}
$c$ is a class function, $c=\sum_\lambda\hat c_\lambda\chi_\lambda$ with $\hat c_\lambda=\langle c,\chi_\lambda\rangle=\frac1{N!}\sum_\sigma c(\sigma)\chi_\lambda(\sigma)$.  The generating function of $c$ against a character is the content polynomial \cite[Ch.~I, \S3 Ex.~4 and \S7]{Mac95}:
\[
\sum_{\sigma\in\SN}x^{c(\sigma)}\chi_\lambda(\sigma)=d_\lambda\prod_{\square\in\lambda}\bigl(x+\cont(\square)\bigr)
\]
(both sides equal $N!\,s_\lambda(1^x)$; for $\lambda=(N)$ it is $x(x+1)\cdots(x+N-1)$).  Differentiating at $x=1$,
\[
\hat c_\lambda=\frac{d_\lambda}{N!}\sum_{\square\in\lambda}\ \prod_{\square'\ne\square}\bigl(1+\cont(\square')\bigr).
\]
A factor $1+\cont(\square')$ vanishes exactly when $\cont(\square')=-1$, i.e.\ for the boxes $(i+1,i)$ on the first subdiagonal.  If $\lambda$ contains two of them, every term of the sum vanishes.  If it contains none, $\lambda=(N)$ and $\hat c_{(N)}=H_N$, the harmonic number (the mean of $c$).  If it contains exactly one, namely $(2,1)$ but not $(3,2)$, then $\lambda_2\ge1$ and $\lambda_3\le1$, i.e.\ $\lambda=(a,b,1^j)$ with $a\ge b\ge1$, and only the term $\square=(2,1)$ survives:
\[
\hat c_\lambda=\frac{d_\lambda}{N!}\prod_{\square\ne(2,1)}\bigl(1+\cont(\square)\bigr)
=\frac{d_\lambda}{N!}\;a!\;(b-1)!\;(-1)^j\,j! ,
\]
since the contents are $0,\dots,a-1$ in the first row, $0,\dots,b-2$ in the second row without its first box, and $-2,\dots,-(j+1)$ down the column.  The hook lengths of $(a,b,1^j)$ are $a+j+1$ at $(1,1)$; $a-c+2$ at $(1,c)$ for $2\le c\le b$; $a-c+1$ at $(1,c)$ for $c>b$; $b+j$ at $(2,1)$; $b-c+1$ at $(2,c)$ for $c\ge2$; and $j,j-1,\dots,1$ down the column.  Their product is
\[
\frac{N!}{d_\lambda}=(a+j+1)\cdot\frac{a!}{(a-b+1)!}\cdot(a-b)!\cdot(b+j)\cdot(b-1)!\cdot j!
=\frac{(a+j+1)(b+j)\,a!\,(b-1)!\,j!}{a-b+1},
\]
and therefore
\begin{equation}\label{eq:chat}
\hat c_\lambda=(-1)^j\,\frac{a-b+1}{(a+j+1)(b+j)}\qquad(\lambda=(a,b,1^j)),\qquad
\hat c_\lambda=0\ \text{ for all other }\lambda\ne(N).
\end{equation}

\subsection*{Step 3: the projections}
$\Pi_\lambda$ is right convolution by a central element of the group algebra, so it commutes with every $R(\rho)$.  By \eqref{eq:step1},
\[
\Pi_\lambda\varepsilon=\bigl(R((u\,v))-I\bigr)\Pi_\lambda c=\hat c_\lambda\bigl(R((u\,v))-I\bigr)\chi_\lambda ,
\]
which vanishes for $\lambda=(N)$ (consistent with $\langle\varepsilon\rangle=0$, i.e.\ $\Pr[u\sim v]=\tfrac12$ under the uniform measure, as $\chi_{(N)}\equiv1$).  Characters are real and satisfy the generalised orthogonality relation $\frac1{N!}\sum_\sigma\chi_\lambda(\sigma)\chi_\lambda(\sigma\rho)=\chi_\lambda(\rho)/d_\lambda$, so
\[
\bigl\|\bigl(R((u\,v))-I\bigr)\chi_\lambda\bigr\|^2
=2\|\chi_\lambda\|^2-2\,\frac{\chi_\lambda((u\,v))}{d_\lambda}=2(1-r_\lambda).
\]
Hence $\nu_\varepsilon(\lambda)=2\hat c_\lambda^{\,2}(1-r_\lambda)$, which with \eqref{eq:chat} is \eqref{eq:main}.

\subsection*{Step 4: the eigenvalue}
In \eqref{eq:r} the first row contributes $\sum_{c=1}^a(c-1)=a(a-1)/2$, the second row $\sum_{c=1}^b(c-2)=b(b-3)/2$, and the column $\sum_{i=1}^j(-1-i)=-j(j+3)/2$.\qed

\begin{remark}
The theorem was checked by brute force over $\SN$ for $N\le7$ (applying $P$ to $\varepsilon$ directly and comparing $\|P^d\varepsilon\|^2$ and $\langle\varepsilon,P^d\varepsilon\rangle$, $d\le4$, with \eqref{eq:corr}; e.g.\ for $N=6$, $d=1$ both give $799/3375$ and $17/45$), and $\sum_\lambda\nu_\varepsilon(\lambda)=1$ was checked to $10^{-12}$ for $N\le4000$.
\end{remark}

\section{Consequences}\label{sec:cons}

\subsection{The two-point correlation and the exact rate}

Write $\Phi(m):=\E[\varepsilon(\sigma_0)\varepsilon(\sigma_m)]=\langle\varepsilon,P^m\varepsilon\rangle$ for the stationary walk, and index the shapes of Theorem~\ref{thm:main} by the number $k=b+j=N-a$ of boxes outside the first row.

\begin{corollary}\label{cor:bound}
For all $N\ge4$ and $m\ge1$,
\[
0\le\|P^d\varepsilon\|^2=\Phi(2d),
\qquad
\bigl|\Phi(m)\bigr|\le\frac{12}{m}+32\Bigl(\frac{4\log(m+1)}{m}+\frac3N\Bigr)^{4}+\frac1{(m+1)^4}.
\]
\end{corollary}

\begin{proof}
Fix $k$ and $b$, so $a=N-k$, $j=k-b$, and $N(N-1)(1-r_\lambda)=k(2N-k-1)+j^2+3j-b^2+3b$.

\emph{Level mass.}  Since $a-b+1\le a+j+1$, $c_\lambda\le1/k$.  For $k\le N-3$, $j^2+3j-b^2+3b\le k^2+3k\le kN$, so $1-r_\lambda\le3k/(N-1)\le6k/N$; for $k\ge N-2$ trivially $1-r_\lambda\le2\le6k/N$.  Thus $\nu_\varepsilon(\lambda)\le12/(kN)$, and as there are at most $k$ shapes with a given $k$, the total mass at level $k$ is at most $12/N$.

\emph{Spectral gap at level $k$.}  Since $b\le k$ and $b\le a=N-k$, $b^2\le k(N-k)$, so $N(N-1)(1-r_\lambda)\ge k(2N-k-1)-k(N-k)=k(N-1)$, i.e.\ $r_\lambda\le1-k/N$ for every shape.  Hence
\[
\sum_{\lambda:\ r_\lambda\ge0}\nu_\varepsilon(\lambda)\,r_\lambda^{\,m}\le\frac{12}N\sum_{k\ge1}e^{-km/N}=\frac{12}{N}\,\frac1{e^{m/N}-1}\le\frac{12}m .
\]

\emph{The bottom of the spectrum.}  For $r_\lambda<0$ we use $|r_\lambda|^m\le(1-\delta)^m$ unless $1+r_\lambda\le\delta$.  Now $N(N-1)(1+r_\lambda)=N(N-1)-j(j+3)+a(a-1)+b(b-3)$, and with $j=N-a-b$ one checks $N(N-1)-j(j+3)=(a+b)(N+j)-N-3j$ and $a(a-1)+b(b-3)\ge-2$; for $a+b\ge3$ this gives $N(N-1)(1+r_\lambda)\ge(a+b-1)N-2$ (the only shape with $a+b=2$ is the sign character $(1^N)$, with $r=-1$ and $\nu_\varepsilon=4/(N^2(N-1)^2)$).  So $1+r_\lambda\le\delta$ forces $a+b\le\delta N+3$; such shapes have $\nu_\varepsilon(\lambda)\le4c_\lambda^2\le64(\delta N+3)^2/N^4$ (as $a+j+1$ and $b+j$ are both $\ge N/2$ once $\delta N+3\le N/2$, which holds whenever the bound is non-trivial), and there are at most $(\delta N+3)^2/2$ of them.  Hence $\nu_\varepsilon([-1,-1+\delta])\le32(\delta+3/N)^4$, and taking $\delta=4\log(m+1)/m$ gives $(1-\delta)^m\le(m+1)^{-4}$.
\end{proof}

The rate $1/m$ is sharp, and the constant is $2$:

\begin{corollary}[continuum limit]\label{cor:cont}
As $N\to\infty$, the spectral measure $\nu_\varepsilon$ converges weakly to the law of $r=1-2(1-\alpha)(1-\beta)$ under the probability density
\begin{equation}\label{eq:density}
\frac{4(\alpha-\beta)^2}{(1-\alpha)(1-\beta)}\,d\alpha\,d\beta
\qquad\text{on }\{\alpha\ge\beta\ge0,\ \alpha+\beta\le1\},
\end{equation}
whose mass within $\delta$ of $r=1$ is $2\delta(1+o(1))$ as $\delta\to0$ and within $\delta$ of $r=-1$ is $O(\delta^4)$.  Consequently
\begin{equation}\label{eq:cont}
\lim_{N\to\infty}\Phi(m)=\int\!\!\int\frac{4(\alpha-\beta)^2}{(1-\alpha)(1-\beta)}\bigl(1-2(1-\alpha)(1-\beta)\bigr)^m\,d\alpha\,d\beta
=\frac{2}{m}+O\Bigl(\frac{\log m}{m^2}\Bigr).
\end{equation}
\end{corollary}

\begin{proof}
Put $a=\alpha N$, $b=\beta N$, $j=(1-\alpha-\beta)N$.  Then $c_\lambda=\frac{(\alpha-\beta)+1/N}{((1-\beta)N+1)(1-\alpha)N}$, $1-r_\lambda=2(1-\alpha)(1-\beta)+O(1/N)$ uniformly, and the shapes in a cell $d\alpha\,d\beta$ number $N^2d\alpha\,d\beta$, so $\nu_\varepsilon$ is a Riemann sum for \eqref{eq:density}; the total mass is $1$ by Theorem~\ref{thm:main}.  Near $r=1$ put $s=1-\alpha$: the region is $\beta\le s$, the density is $4/s+O(1)$, and $1-r=2s(1-\beta)=2s+O(s^2)$, so the mass with $1-r\le\delta$ is $\int_0^{\delta/2}4\,ds\,(1+O(\delta))=2\delta+O(\delta^2)$; integrating $4(1-2s)^m$ over $s\le\tfrac12$ gives $2/(m+1)$ and the corrections are of relative order $\log m/m$.  Near $r=-1$ both $\alpha$ and $\beta$ are $O(\delta)$ and the density is $O(\delta^2)$ on a set of area $O(\delta^2)$.
\end{proof}

\begin{remark}
The mass $2\delta$ at the top of the spectrum has a transparent probabilistic meaning.  Under the uniform measure the pair of cycle lengths $(|C_u|,|C_v|)$ is well understood: the smaller of the two relevant masses (the shorter arc between $u$ and $v$ when $u\sim v$, the smaller of the two cycles otherwise) has density $4$ at $0^+$ on the scale of $N$, and a pair whose smaller mass is $g$ keeps its status for $\asymp N/g$ steps, since no transposition touches the small part.  Thus the slow modes of $\varepsilon$ are the ``trapped'' pairs, and the linear spectral density is the linear density of the gap.  Theorem~\ref{thm:main} shows that this trapping is the \emph{only} slow mechanism and identifies the constant.
\end{remark}

\subsection{One step}

\begin{corollary}\label{cor:onestep}
Let $\sigma$ be uniform on $\SN$ and $\tau$ an independent uniform transposition.  Then
\[
\Pr\bigl[u\sim_\sigma v\ \text{ and }\ u\not\sim_{\sigma\tau}v,\ \text{ or vice versa}\bigr]
=\frac{(N+1)(N+2)}{6N(N-1)},
\qquad
\Phi(1)=1-\frac{(N+1)(N+2)}{3N(N-1)} .
\]
\end{corollary}

\begin{proof}
$\Phi(1)=1-2q$ where $q$ is the probability in question, and $\Phi(1)=\sum_\lambda\nu_\varepsilon(\lambda)r_\lambda$ by \eqref{eq:corr}.  The sum is a rational function of $N$ which can be evaluated in closed form, but it is quicker to compute $q$ directly: conditionally on $u\sim v$ with common cycle of length $\ell$ (probability $(\ell-1)/(N(N-1))$ for each $\ell\ge2$, the $u\to v$ distance $x$ being uniform on $\{1,\dots,\ell-1\}$), $\tau$ separates the pair iff its two points lie on different arcs, probability $2x(\ell-x)/(N(N-1))$; conditionally on $u\not\sim v$ with cycle lengths $(p,q')$ (probability $\frac1N\cdot\frac{N-p}{N-1}\cdot\frac1{N-p}$ for each admissible pair, by choosing $u$'s cycle and then $v$'s), $\tau$ merges them with probability $2pq'/(N(N-1))$.  Summing the two finite sums gives $q\cdot\binom N2=\frac{(N+1)(N+2)}{12}$.
\end{proof}

So one random transposition changes the status of a fixed pair with probability $\tfrac16+\tfrac{2}{3N}+O(N^{-2})$, and $\Phi(1)\to\tfrac23$.

\section{Generalisations}\label{sec:gen}

\subsection{Other conjugation-invariant walks}
Nothing in Steps 1--3 refers to the transposition walk except the final evaluation $\chi_\lambda((u\,v))/d_\lambda=r_\lambda$, which comes from the \emph{definition} of $\varepsilon$, not from the walk.  Hence for any random walk $\sigma_{t+1}=\sigma_t\rho_{t+1}$ with $\rho$ drawn from a conjugation-invariant law $\mu$ on $\SN$ (random $k$-cycles, random involutions, the $(N-k)$-cycle walk, \dots), whose eigenvalues are $r^\mu_\lambda=\sum_\rho\mu(\rho)\chi_\lambda(\rho)/d_\lambda$,
\[
\E\bigl[\varepsilon(\sigma_0)\varepsilon(\sigma_m)\bigr]=\sum_{\lambda=(a,b,1^j)}2c_\lambda^{\,2}(1-r_\lambda)\,\bigl(r^\mu_\lambda\bigr)^m ,
\]
with $c_\lambda$ and $r_\lambda$ as in Theorem~\ref{thm:main}.  The spectral measure is intrinsic to the function $\varepsilon$; only the moments change.

\subsection{Non-stationary starts}
For a fixed start $\sigma_0$, $\Pi_\lambda\varepsilon=\hat c_\lambda(R((u\,v))-I)\chi_\lambda$ gives
\begin{equation}\label{eq:fixed}
\E\bigl[\varepsilon(\sigma_0\tau_1\cdots\tau_m)\bigr]=(P^m\varepsilon)(\sigma_0)
=\sum_{\lambda=(a,b,1^j)}r_\lambda^{\,m}\,\hat c_\lambda\,\bigl[\chi_\lambda(\sigma_0(u\,v))-\chi_\lambda(\sigma_0)\bigr],
\end{equation}
with $\hat c_\lambda$ carrying the sign $(-1)^j$ of \eqref{eq:chat}.  At $m=0$ this reads $\sum_\lambda\hat c_\lambda[\chi_\lambda(\sigma_0(u\,v))-\chi_\lambda(\sigma_0)]=\varepsilon(\sigma_0)$, which for $\sigma_0=\mathrm{id}$ is the identity $\sum_{\lambda=(a,b,1^j)}\hat c_\lambda d_\lambda(1-r_\lambda)=1$.  When $\sigma_0$ is an $N$-cycle, $\chi_\lambda(\sigma_0)$ vanishes unless $\lambda$ is a hook, and $\sigma_0(u\,v)$ has cycle type $(\ell,N-\ell)$, where the characters of the shapes $(a,b,1^j)$ are given by the Murnaghan--Nakayama rule in terms of two rim hooks; \eqref{eq:fixed} then becomes an explicit finite sum for the probability that two points at distance $\ell$ on an $N$-cycle are still (or again) in a common cycle after $m$ random transpositions.

\subsection{The structure behind the formula}
The identity \eqref{eq:step1} says that the same-cycle indicator is a first difference of the cycle count, and the cycle count is supported on the shapes $(a,b,1^j)$ because its generating function is the content polynomial, whose derivative at $x=1$ survives only one subdiagonal box.  More generally, for any $g\in\SN$ the function $\sigma\mapsto c(\sigma g)-c(\sigma)$ has spectral measure $\hat c_\lambda^{\,2}\,\|(R(g)-I)\chi_\lambda\|^2=2\hat c_\lambda^{\,2}(1-\chi_\lambda(g)/d_\lambda)$, supported on the same shapes.  Functions such as $\mathbf 1\{u_1\sim u_2\sim u_3\}$ are not of this form, and we do not know whether their spectral measures are equally explicit.

\begin{remark}[the Latin square problem]\label{rem:latin}
Let $L$ be a Latin square of order $n$, let $j,j'$ be two columns, and let $\pi$ be the permutation of the rows sending $r$ to the row $r'$ with $L(r',j')=L(r,j)$ (the ``column cycles'' of the pair).  Whether two given rows lie in a common cycle of $\pi$ is a bit that appears in the switching identities of Cavenagh, Greenhill and Wanless for the distribution of the second row of a random Latin square \cite{CGW08}, and the re-randomisation method of Kwan, Petrova and Sawhney \cite{KPS25} randomises it by switching a pool of disjoint intercalates through the columns $j,j'$, each of which multiplies $\pi$ by a transposition of rows.  If the base permutation is uniform and the pool consists of $k$ uniform disjoint transpositions, the bit's conditional bias given the pool, $\mathrm{bias}=\Pr_S[\text{bit}]-\tfrac12$ over uniform subsets $S$ of the pool, satisfies $4\E[\mathrm{bias}^2]=\E_{M}\Phi(M)+O(k^2/n)$ with $M\sim\mathrm{Bin}(k,\tfrac12)$, so that $\E|\mathrm{bias}|=O(k^{-1/2})$ by Corollary~\ref{cor:bound}; this is the estimate that the present computation was made for.
\end{remark}

\begin{thebibliography}{9}
\bibitem{BD06} N.~Berestycki and R.~Durrett, A phase transition in the random transposition random walk, \emph{Probab.\ Theory Related Fields} 136 (2006), 203--233.
\bibitem{BSZ11} N.~Berestycki, O.~Schramm and O.~Zeitouni, Mixing times for random $k$-cycles and coalescence-fragmentation chains, \emph{Ann.\ Probab.} 39 (2011), 1815--1843.
\bibitem{CGW08} N.~J.~Cavenagh, C.~Greenhill and I.~M.~Wanless, The cycle structure of two rows in a random Latin square, \emph{Random Structures Algorithms} 33 (2008), 286--309.
\bibitem{DS81} P.~Diaconis and M.~Shahshahani, Generating a random permutation with random transpositions, \emph{Z.\ Wahrsch.\ Verw.\ Gebiete} 57 (1981), 159--179.
\bibitem{KPS25} M.~Kwan, A.~Petrova and M.~Sawhney, Random Latin squares: the parity of the second row, arXiv:2509.13125 (2025).
\bibitem{Mac95} I.~G.~Macdonald, \emph{Symmetric Functions and Hall Polynomials}, 2nd ed., Oxford University Press, 1995.
\bibitem{Sch05} O.~Schramm, Compositions of random transpositions, \emph{Israel J.\ Math.} 147 (2005), 221--243.
\end{thebibliography}

\end{document}
